\chapter{How to Write a Book: ベストセラー作家による13のステップ}

この講義は、意図的に組まれた制約と解放の連鎖として展開する。私たちはまず、講師がそうするように、 gimmick によって本がすばやく作れるという安易な internet の約束に抗するところから始める。ついで、一歩ずつ、作業条件から規模へ、規模から構造へ、構造から reader の前進感へ、そして最後には、drafting の自由から ruthless な revision へと進んでいく。そこから現れるのは、tips の寄せ集めではなく、長い manuscript がどのように built されるかについての disciplined な model である。

\section{まず quality: space、tools、そして excuses の拒否}

講義は polemical に始まる。speed を約束するもの、instant success の formulas、bestseller を生み出すための compressed な recipes は世の中にいくらでもある。講師は、その atmosphere をただちに reject する。彼は bestseller を promise しないし、speed も promise しない。promise するのは thirteen foundational steps である。支配的な contrast は、簡潔には
\begin{equation}
\text{speed} \neq \text{goal}, \qquad \text{quality} = \text{goal}.
\end{equation}
と書ける。この relation が、その後に続く everything の tone を set する。もし quality が goal なのだとすれば、私たちの first concern は acceleration ではなく、sustained work が possible になる conditions なのである。

だからこそ講義は、これほど concretely に始まる。私たちは writing space を establish する。ついで writing tools を assemble する。これら最初の steps は glamorous ではない。だが、beginner の favorite alibi を eliminate するために置かれている。講師は explicit である。書く場所がない、と言ってはならない。かつては、chairs の上に plank を渡したものが desk の代わりになった。restaurant でもよい。car でもよい。noisy な newsroom ですら一時は役に立った。point は、ideal conditions が irrelevant だということではない。point は、work は merely workable な conditions のもとで begin する、ということだ。

同じ seriousness は tools にも apply される。elaborate な apparatus は要らない。だが、数分ごとに concentration を break してしまわない enough な stability は必要である。こちらを fight しない computer、back や neck を punish しない chair、そして otherwise room を行ったり来たりして search することになる ordinary な materials。これらはすべて同じ logic に属している。どんな large project であれ move し始める前に、まず friction を reduce するのである。

\subsection{Question \& Answer}

\paragraph{Question.} begin する前に ideal conditions は必要なのか。

\paragraph{Answer.} No である。講義は、ideal conditions と usable conditions を distinguish する。privacy、silence、better equipment はたしかに help する。だが、それらの absence は delay の excuse にはならない。私たちに必要なのは、recur できる place と setup である。operative word は perfect ではなく established なのである。

\section{problem の scale を break し、idea を test する}

local な conditions が fixed されると、講義は次の obstacle を raise する。本とは sheer な mass だけで人を intimidate するに十分なほど large である。講師は、これに inspiration では answer しない。think する scale を changing することで answer する。本は、five-hundred-page object として attack されるのではない。それは units に reduce される。
\begin{equation}
\text{book} \to \text{chapters} \to \text{paragraphs} \to \text{sentences}.
\end{equation}
この reduction こそが、elephant analogy の practical content である。私たちは whole を一度に consume しない。次の一片を consume するのである。

同じ point は、より schematically にも書ける。
\begin{equation}
\text{colossal task} \Rightarrow \text{small actionable units}.
\end{equation}
講義がここに linger するのは、これが最初の major な psychological reversal だからである。本それ自体はなお large なままだが、writer の next action は small になる。

だがこの point で、講師は easy な conclusion を refuse する。すべての project が perseverance によって rescue されるに値するわけではない。ある ideas は、いま occupy するよう求められている form に対して too small なのである。そこで Step 4: settle on a big idea となる。この distinction は vague ではない。small concept は blog post や article に belong する。big concept は book に belong する。
\begin{equation}
\text{small concept} \Rightarrow \text{blog/article}, \qquad
\text{big concept} \Rightarrow \text{book}.
\end{equation}
講師は、abstract に語るのではなく examples によってこれを vivid にする。nonfiction については \textit{How to Win Friends and Influence People} を invoke し、fiction については \textit{Harry Potter} を invoke する。point は imitation ではない。scale である。

\subsection{Question \& Answer}

\paragraph{Question.} idea に本を支える enough な legs があるかどうかを、どう知るのか。

\paragraph{Answer.} 講義は pragmatic な test を与える。idea を another person に tell してみるのである。retelling のたびにそれが bigger になっていくなら、おそらく required な scope を持っている。繰り返し同じ wall に hit するなら、trouble は stamina ではない。idea 自体が undersized なのである。

\section{outline、fiction と nonfiction、そして marathon of the middle}

Step 5 で講義は最初の structural summit に reach する。whiteboard が matter するのは、一連の practical rules が、shape の theory になるその瞬間を capture しているからである。

\begin{figure}[h]
\centering
\includegraphics[width=0.88\textwidth]{lecture_05_figure_03.png}
\caption{Step 5 の whiteboard outline: fiction、nonfiction、the middle、そして outline の role。}
\end{figure}

整えた typeset form では、board は次のように read される。
\begin{align}
5.\ \text{OUTLINE} \notag\\
\text{Fiction vs.\ Non-fiction}\qquad (\text{King})\qquad (\text{Agents/editors}) \notag\\
\text{The Marathon of the Middle} \notag\\
\text{Same structure} \notag\\
\text{Outline serves } \underline{\text{you}}. \notag
\end{align}

講師は、audience が feel するであろう tension から始める。もし some fiction writers が detailed plan ではなく discovery によって work する pantsers なのだとしたら、なぜ outline に insist するのか。その answer は careful である。outlining の degree は vary してよい。だが direction は disappear できない。Stephen King は one pole を stand している。interesting characters が difficult situations に placed され、writing を使って何が happen するかを discover する pole である。だが even there、some sense of direction はなお necessary である。nonfiction は opposite pole に lie する。そこでは writer は、何が said われ、どう said われ、sources は何で、argument はどこへ going なのかを know したい agents や editors に answerable である。

だからこそ、board の next claim が so much matter するのである。 \textit{same structure}。surface resources は genres を across して differ する。だが underlying motion は differ しない。講義は、この shared law に simple な form を与える。
\begin{equation}
\text{setup} \to \text{payoff}.
\end{equation}
これが、この chapter の real mathematical spine である。fiction はしばしば scenes、dialogue、reversal、conflict を通してこの pattern を realize する。nonfiction は、problem、method、explanation、result を通して realize するかもしれない。だが in both cases、reader は planted expectation と later fulfillment によって carry されるのである。

\begin{definition}
The Marathon of the Middle とは、manuscript の portion、通常は halfway point と three-quarter point の somewhere の between に位置する部分であり、そこでは writer が length の drag と、pad したり abandon したり、merely endure したりしたい temptation を feel する。
\end{definition}

compact な editorial shorthand としては、
\begin{equation}
x \in \left[\tfrac12,\tfrac34\right] \Rightarrow \text{heightened risk of structural fatigue},
\end{equation}
と書ける。ただし variable \(x\) は私たち自身の summary であって、lecture で使われた symbol ではないことを remember しておくべきである。

ここでの講師の strongest claim は、middle を merely survive してはならない、ということである。そこで thrive しなければならない。middle が us に boring に seem するなら、reader は already asleep なのだ。この principle は nonfiction にも exactly as much apply する。model ship の building に関する nonfiction book は、merely instructions を enumerate してはならない。eventual solution が earned に feel されるよう、difficulty を setup しなければならないのだ。

\subsection{Question \& Answer}

\paragraph{Question.} もし私たちが pantsers なら、なぜそれでも outline を need するのか。

\paragraph{Answer.} discovery は aimlessness と同じではないからである。fiction writer は lightly に outline してよく、nonfiction writer は heavily に outline してよい。だがいずれの場合も、outline は、otherwise sprawl してしまうだろう long form に direction を与える。lecture の final corrective も equally important である。outline は instrument である。それは book に serve し、book は writer に serve する。developing manuscript が original scheme より better だと prove したなら、book をそれに subordinate するのではなく、outline のほうを change すればよいのだ。

\section{schedule、deadlines、そして control variables としての research}

structure が established されると、lecture は execution に return する。Step 6 は、write したいという will を arithmetic へ turn する。writer は manuscript の total size を estimate し、available な days を estimate し、divide する。
\begin{equation}
\text{daily quota} = \frac{\text{target pages or words}}{\text{days available}}.
\end{equation}
この formula の significance は numerical elegance にはない。moral clarity にある。deadline が real になるのは、それが daily obligation に convert されたときだけなのだ。

lecture は equally practical に adjustment についても語る。私たちは initially に ten pages per day と estimate して begin するかもしれない。だが later に、four or five のほうが reality に closer だと discover するかもしれない。その場合 schedule は revise されうる。だが once revised and fixed されたなら、それは sacred になる。mistake は、writing のための time は accidentally に found されるだろうと believe することだ。lecture は逆を言う。day の中の何か else が sacrifice されねばならないのである。

\paragraph{Worked example.} speaker は、自分自身の schedule から concrete な case を give する。
\begin{align}
70{,}000 - 67{,}500 &= 2{,}500,\\
\text{tomorrow's quota} &= 2{,}500\ \text{words}.
\end{align}
calculation 自体は elementary である。だが point は elementary ではない。vague ambition を sharp demand へ transform することにある。

講師はさらに sobering な publishing heuristic を add する。
\begin{equation}
\text{deadline keepers} \approx \frac{1}{100}.
\end{equation}
これは theorem ではない。それは、seriousness の standard を raise するための professional estimate である。deadline を守ることは、それだけですでに、most writers から oneself を distinguish することなのだ。

Step 7 はそれから discussion を、schedule から knowledge へ broaden する。research は nonfiction で obviously necessary である。だが many writers は、fiction におけるその importance を underestimate している、と講師は insist する。history、aircraft、weaponry における missed detail は、その subject を know している readers に immediately noticed されるだろう。lecture を govern する metaphor は exact である。
\begin{equation}
\text{research} \neq \text{main course}, \qquad
\text{research} = \text{seasoning}.
\end{equation}
この metaphor は hierarchy を preserve する。research は specificity と credibility を supply するのであって、story に replace するのではない。講師は自身の work から concrete な example を offer する。2000 B.C.\ から present に span する world-history chart を使って、births、deaths、overlaps を accurately に track しているのだ、と。research は narrative line を support するために exist しているのであって、それを swamp するためにあるのではない。

\subsection{Question \& Answer}

\paragraph{Question.} book を takeover し始める前に、どれくらいの research が enough なのか。

\paragraph{Answer.} work を specific にし、believable にし、technically secure にするに enough でありながら、reader が writer の note cards を admire するよう forced されるほどではない程度である。research がその job を done したと言えるのは、page の subject になることなく、page を sharpen したときである。

\section{reader-first な openers と tension の logic}

structure、schedule、research が in place すると、lecture は decisively に reader へ turn する。Step 8 は opening line から begin する。point は、first line が magical だということではない。tone を fix し、pressure を establish するということだ。講師はいくつかの categories of opener を name する。
\begin{equation}
\{\text{surprising},\ \text{dramatic},\ \text{philosophical},\ \text{poetic}\}.
\end{equation}
そして examples を与える。\textit{1984} における Orwell の opening は、world を immediately destabilize する surprise の illustration として used される。Elizabeth Gilbert の nonfiction example も、sharp で concrete な action によって同じ end に向かう。 \textit{Paradise} における Toni Morrison の famous な opening は dramatic statement を stand している。lecture はしかし、example の level に long く stay しない。first line という local issue を、global rule へ convert するのである。

その rule が、reader-first filter である。
\begin{equation}
\text{decision} \xrightarrow{\text{reader first}} \{\text{keep},\ \text{revise},\ \text{cut}\}.
\end{equation}
講師はこれを directly に言う。manuscript における every decision は、それが reader に何を do するのかという question を pass しなければならない。writer first ではない。editor first でもない。critic first でもない。reader first である。これは talk における strongest conceptual turns の一つである。opener が important なのは ornament としてではなく、それがあとの all later decisions がどう judged されるべきかを teach してくれるからである。

Step 9 は、その principle を opening line から whole manuscript の life へ extend する。何が reader に page を turn させ続けるのか。lecture の answer は pressure である。
\begin{equation}
\text{conflict/tension} \Rightarrow \text{page-turning momentum}.
\end{equation}
fiction において、その explanation は vivid である。pleasant な agreement は life においては nice かもしれない。だが page の上では、それは often die する。scene が alive になるのは、some buried strain が appear し、一人の character が unexpected thing を say し、surface の conversation がもはや自分自身を explain しなくなるときである。

しかし nonfiction も exempt ではない。講師は explicitly に、there で tension がどう appear するのかを ask する。その answer は theatrical ではなく structural である。
\begin{equation}
\text{promise} \to \text{delivery}.
\end{equation}
strong な nonfiction book は、early pages において、reader が end までに何を possess することになるのかを promise する。そして later に、それを exactly に deliver しなければならない。この context における tension は battle ではない。deferred fulfillment である。

\subsection{Question \& Answer}

\paragraph{Question.} literal villains や battles や melodrama がなくても、nonfiction に tension を create するにはどうすればよいか。

\paragraph{Answer.} reader に wait する何かを与えることである。difficulty、不可能、あるいは desired transformation を pose し、そのあと manuscript を organize して、later chapters が earlier chapters の promised したものを fulfill するようにするのである。nonfiction において、tension はしばしば、announcement と arrival のあいだの disciplined な spacing にすぎない。

\section{drafting の freedom、return of the middle、そして resounding な ending}

Step 10 は temporal register を change する。ここまでは mostly structure を outside から describe してきた。now it asks、drafting という act の during に writer が何を permit すべきかを。answer は severe であると同時に liberating でもある。turn off the internal editor、である。
\begin{equation}
\text{draft with editor off} \to \text{return later to edit} \to \text{cut to add power}.
\end{equation}
この sequence が matter する。first draft は、every word を litigate する place ではない。講師はこれを personal terms で言う。inner critic は shoulder に sit して、everything に object する。drafting の during には、その critic は silence されねばならない。first draft の purpose は、material を get down することなのである。

\subsection{Question \& Answer}

\paragraph{Question.} internal editor は、いつ silent であるべきで、いつ return すべきなのか。

\paragraph{Answer.} draft が being produced のあいだは silent であるべきである。early criticism は continuity を destroy するからである。それは only later に、preferably some time has passed したあとで、writer が manuscript を、still struggling to come into existence しているものではなく、shaped されるべきものとして look できるようになったときに return すべきなのである。

この point で lecture は、もっとも important な recaps の一つを perform する。Step 11 は explicitly に Marathon of the Middle へ return する。これは repetition for its own sake ではない。earlier の structural principle が recurring lived crisis として reintroduced される moment なのである。講師は plainly に、ここが彼自身も quit したくなる place だと言う。danger は padding である。fiction における extra scenes、nonfiction における extra points、pressure を increase させることなく length を prolong する material である。

answer は again、endure ではなく thrive である。middle は real な setup と payoff を continue to carry しなければならない。講師自身の current work in progress では、middle を alive に keep する one strategy は、2000 B.C.、present day、Vietnam のあいだを alternate することだ。manuscript が single band of exposition へ flatten する代わりに、自らの pressure を renew し続けるためである。

Step 12 は、それから closure へ turn する。ending は fizzle してはならない。それは theater における curtain が down するような firmness をもって land しなければならない。if several endings are possible であるなら、lecture の rule は clear である。
\begin{equation}
\text{candidate endings} \to \text{choose the most emotional}.
\end{equation}
これは sentimentality を mean しない。読者は自分を move したものを remember するのであり、したがって book の last movement は、それまで all along prepared されてきた force を gather しなければならないのだ。

\section{ferocious な self-editing と professional judgment}

lecture の final movement は、その standard を last time でもう一段 tighten する。私たちは ferocious self-editors になれと told される。この word は mystical ではない。aggressive であれ、ということだ。私たちは cut する。polish する。manuscript を sing させる。もしこの work を自分たち自身でしないなら、それ以前の everything は partly wasted されているのである。

講師はついで、manuscripts が actual に read される professional pressure を introduce する。agents や editors には、too much to read があり、too little time しかない。very often、彼らは two minutes の within に、あるいは one or two pages の within に、その manuscript が further attention に deserve するかどうかを tell できる。logic は harsh だが simple である。competition は vast だから、opening pages はすでに sharpened された状態で arrive しなければならない。したがって self-editing は decorative last step ではない。book が apology なしに world に meet する ready になるための act なのである。

最後に lecture は、自分自身の beginning に circle back する。それは shortcuts と speed-hype を rejecting するところから始まった。そしてそれらの place に何が stand するのかを示して終わる。working conditions、big idea、outline、schedule、research、reader-first judgment、tension、free draft、thriving する middle、landing する ending、そして manuscript が its full force を carry するまで cut するに十分 ruthless な editor である。

\section{まとめ}

lecture は one step at a time で advance する。だがその force は、steps が accumulate していく way に lie している。まず excuses を remove し、work のための conditions を establish する。problem の scale を reduce し、ついで concept が本を deserve するほど large かどうかを test する。scheme の slave としてではなく、long form を通して direction を need とする writers として outline する。ambition を quotas へ convert し、research を display ではなく seasoning として use し、every local decision を reader の experience によって judge する。draft を edit する前に write し、それを alive にすることによって middle を survive し、もっとも deeply に move する ending を choose し、最後に ferocity をもって revise する。この lecture の seriousness は、一つの simple truth から来ている。本は wish されることで existence に入るのではない。それは built され、sustained され、sharpened されるのである。